\documentclass[11pt]{article}
\usepackage[T1]{fontenc}
\usepackage{lmodern,microtype}
\usepackage{amsmath,amssymb,enumitem,booktabs}
\usepackage[letterpaper,margin=1in]{geometry}
\usepackage[colorlinks=true,urlcolor=blue,linkcolor=blue]{hyperref}
\title{On a metaplectic Jacquet--Langlands transfer for division algebras of prime degree\\
(book) --- Revision Log for Draft 1}
\author{Prepared by Claude (Anthropic) for David Joyner}
\date{November 8, 2026}
\begin{document}
\maketitle

\section*{Decisions (David Joyner, November 8, 2026)}
\begin{enumerate}[leftmargin=*]
\item Combine Parts~I--IV and the companion paper on the rank-two wild
place into one book, so that cross-references and citations work.
\item Freeze the separate papers: their current drafts (Part~I
Draft~11, Part~II Draft~10, Part~III Draft~10, Part~IV Draft~6,
companion Draft~8) stay in the repository as history, and the book is
the working document from now on, with its own drafts and logs.
\item Title: \emph{On a metaplectic Jacquet--Langlands transfer for
division algebras of prime degree}.  Authors on the title page:
David Joyner, Claude (Anthropic), ChatGPT (OpenAI); each chapter's
source and author order are listed in the Preface.
\end{enumerate}

\section*{Structure}
\begin{tabular}{@{}ll@{}}\toprule
Front matter & Preface (sources, conversion, contributions, the five acknowledgments)\\
1 & Introduction (new; setting, chapters, running example, open problems)\\
2 & Literature: Part~I's review, then the related-work sections of Part~II,\\
  & Part~III and the companion paper\\
3 & Notation and glossary: the shared glossary, then the notation sections\\
  & of Parts~I, II, III and IV\\
4--8 & Parts~I, II, III, IV and the companion paper\\
9 & Summary of results, one section per chapter (Part~IV: its abstract)\\
Appendix A & Tables and verification scripts of the companion paper\\
Bibliography & merged, 104 entries\\
\bottomrule\end{tabular}

\noindent Each chapter opens with the source's abstract (``Overview'')
and a line naming the source draft.  Numbering is by chapter and
section (Theorem~4.4.5).  The layout follows David's choices for Part~I:
literature and notation at the front, summary at the end.

\section*{No mathematical change}
No statement, proof, hypothesis or status label was changed.  Counting
environments in the five sources and in the book gives identical
totals: 26 theorems, 50 lemmas, 40 propositions, 14 corollaries,
16 definitions, 66 remarks, 5 examples, 128 proofs, 124 equations,
6 align blocks, 56 enumerations, 7 tables, 4 diagrams.  Text removed:
the five titles, abstracts (moved to the chapter overviews), the
\verb|\input| of the glossary in Parts~I--III (the glossary appears once,
in Chapter~3), the lead-in paragraphs of the three per-Part ``Notation and
glossary'' sections (replaced by one lead-in for Chapter~3), and the
five bibliographies (merged).  The acknowledgments moved to the Preface.
The build has no errors, no undefined references or citations, and no
multiply defined labels.

\section*{Conversion rules}
\begin{enumerate}[leftmargin=*]
\item \textbf{Labels.}  Every label of a source is prefixed by its
chapter tag (\texttt{I:}, \texttt{II:}, \texttt{III:}, \texttt{IV:},
\texttt{W:}); the shared glossary's labels (\texttt{gl:}) are unchanged.
Four headings that had no label but are cited from other Parts were
given one: Part~I \S1.1 and \S1.4, Part~III \S2, Part~IV \S2.1.
\item \textbf{Cross-Part references.}  Typed references such as
``Part~I, Lemma~2.2'' were mapped, by number, to labels using the
compiled label tables of the drafts they were written against: Parts~II
and~III and the glossary cite Part~I Draft~11; Part~IV, the companion
paper and the running-example note cite Part~I Draft~8 (= Draft~10)
numbering; Parts~II--IV have one numbering each.  Theorem-like items are
rendered with \verb|\cref|, which prints the item's actual type;
sections with \verb|\S\ref|.  ``Part~X'' without a number became
``Chapter~\ref{...}''.  43 references typed without a ``Part'' prefix
(the comparison table, diagram and restatements (I1)--(I11) of Part~II;
section pointers in Part~III's lead-ins; one in Part~I; five in the
companion paper's list of dependencies; one ``Appendix~B'' in Part~IV)
were mapped by hand; each is listed in \texttt{reports/assemble\_fixes.txt}.
\item \textbf{Name mismatches.}  The companion paper cites ``Part~I,
Definition~2.2'' twice (in its section on the cover in the Kazhdan--Patterson model,
now \S8.11, and in the dependency list of its main theorem).  In Part~I's Draft~8 numbering, item~2.2 is Lemma~2.2 (scalar
centrality); Definition~2.1 defines the local covers and \S2.2 is the
ambient cover.  The book renders the reference as \verb|\cref| of item~2.2,
so it now reads ``Lemma~4.2.2''.  David: please check which item was
meant (probably Definition~2.1, now Definition~4.2.1).
\item \textbf{Citations.}  Bracketed citations (Parts~I, II and the
companion paper) were converted to \verb|\cite| only when every key in
the bracket is a key of that paper's bibliography, so that labels such as
[P] and commutators such as $[Q,Q]$ were left alone.  Citations of the
Parts themselves (\texttt{P1}, \texttt{JoynerI}, \ldots,
\texttt{RankTwo0}) became chapter references.
\item \textbf{Macros.}  Each chapter keeps its source's macros,
redefined at the start of the chapter inside a group, so different
definitions of the same name (for example \verb|\GDb|, \verb|\GL|) do not
interact.  Theorem environments are defined once.
\end{enumerate}

\section*{Bibliography}
One key per work.  31 works appeared under several keys and were merged:
BD, Bad, Ber, Coh, DKV, FK, Fli80, GG18, HC70, J, JC, JL, KP1, KP2, Kaz84,
Klo86a, Kot84, Kub67, Kub69, Li2, Li4, Moo, PRa88, Rog83, Sav88, Ser,
Wal80, Wal91, Wei, Wei18, Z (sources listed in
\texttt{reports/bib\_report.json}).  Two key collisions were separated:
[BZ] is Bernstein--Zelevinsky 1976 (Russian Math.\ Surveys) in Part~I but
1977 (Ann.\ ENS) in Part~II, now [BZ76] and [BZ77]; [Gao] is Gao's
$\mathrm{Mp}_4$ preprint in Part~I but his 2018 Amer.\ J.\ Math.\ paper in
the companion paper, now [Gao12] and [Gao18].  Where entries differed,
Part~I's (checked) entry was used, except DKV (the companion paper's
entry has pages) and Cohen (Part~III's has the journal data).  Weil's
\emph{Basic Number Theory} is cited in the 2nd edition by Part~I and the
3rd by the companion paper; the entry names both.  Moore's IH\'ES paper
now has the correct pages 5--70 (the companion paper had 157--222).
Henniart's paper, in Part~I's bibliography but never cited, was dropped.

\section*{Known limitations of Draft 1}
\begin{enumerate}[leftmargin=*]
\item Each chapter keeps its own introduction, so the setting, the
labels and the running example are repeated; prose still says ``this
paper'' and, in places, ``here cited as Draft~8''.  An editing pass
should remove the repetition and update the wording.
\item Parts~IV and the companion paper had not had the exposition pass
(summary, related work, lead-ins) given to Parts~I--III; their
existing summary and literature sections were used.
\item The chapter titles are those of the source papers; Chapter~8's
title is shortened.
\end{enumerate}

\section*{Files}
\texttt{book/mjl\_book\_draft1.tex} and \texttt{.pdf};
\texttt{book/chapters/*.tex} (generated);
\texttt{book/tools/} (the scripts that generated them, with a README);
\texttt{book/reports/} (conversion report, hand fixes, bibliography
report, source bibliographies); this log.  The frozen papers are
unchanged.  README updated.
\end{document}
